\section*{Acknowledgements}

It takes a \sout{village} observatory to raise a PhD student; it would be remiss of me not to acknowledge that dozens of people have helped me in my journey in completing this body of work, whether it be in a personal or in a professional capacity, so this is a non-exhaustive list of acknowledgements.

Firstly, I want to thank my supervisor, Evan. Four years ago, when he asked me to do this project, I was quite unsure about doing a PhD (or anything in general for that matter). I can safely say that any doubts I had about doing a PhD were quickly quieted. The last four years have been the best of my life, largely thanks to a supervisor who nurtured my crazy ideas and let me fly around the globe pursuing them.  It’s been one of my great privileges to learn from Evan, and hopefully we’ll get up to some more crazy stuff in the future. Slán go fóill fear fiáin. 

Additionally, I want to thank my co-supervisor, Vishal, with whom I work closely in trying to answer the hardest question I could think of: Is there intelligent life in the universe? I have great memories of our meetings in Campbell, all the way to Hat Creek. Vishal’s thoughts and guidance on all things SETI were indispensable in completing this PhD. 

I also want to thank Joe McCauley. Joe made my thesis projects possible from a hardware point of view, from the long drives to Birr, to the many hours spent fixing and discussing servers along with many interesting things and a couple pints in-between it all. Without Joe, the I-LOFAR station would have fallen into the abyss long ago, of that I am sure. 

I have benefited greatly from discussions with Aaron Golden throughout my PhD. I want to acknowledge the significant time and thought Aaron put into conversations about life, science and careers alike, even though at times it took a couple of dictionaries to parse what he was saying. I will remember the trials and tribulations we faced on our journey back home from India for the rest of my life. 

I want to specifically highlight the help David McKenna has given me over the past few years. Despite David submitting his PhD and starting a new, intensive job throughout my tenure as a PhD student, he always took time to help me with silly questions as I tried to find my feet in the wacky world of pulsars and radio astronomy. David’s knowledge of the field is incomprehensible, and for that I am grateful. 

To my PhD brother Sai Susarla, thank you for showing me all the tricks with timing, solar wind and Indian food. I had a blast every time I came to work and hang out in Galway. I suspect we have a couple more papers to write together before its all said and done. 

Thanks to Harry Qiu, for many things, but namely helping me interface with Parkes, simulating pulses and giving recommendations for café's in Hong Kong. 

The radio group at Trinity, Lety, Conor, Charlie, Jessica and Karen. Thank you all for your help, company and putting up with my meanderings about the transient sky. It was a pleasure to work with each of you and I have no doubt you'll all be very successful in your respective works. Being part of the radio rodeo is for sure one of the highlights of this experience. 

A big thanks to all the people who I shared SNIAM's penthouse over the last four years, especially Sorcha and Kevin who I started my PhD with. The office was always a great environment to spend my days working, laughing and goofing off. I will miss seeing you all everyday. 

I want to thank Dave DeBoer and the whole LFT3 team, working on the development of this project has been immensely satisfying and is accumulation of everything I wanted to do as an astronomer, I am glad that we do things not because they are easy but because they are hard. 

A substantial part of my PhD was spent at the University of California, Berkeley working with Breakthrough Listen (BL). I want to thank Andrew Siemion BL's director for providing the resources to make the expansion of SETI efforts to Ireland and this PhD possible. I also want to thank the Breakthrough Foundations chairmen Pete Warden for the many conversations and advice given on everything to do with space industry and agencies. Along with Steve Croft for all his efforts in running the undergraduate research program, a program I participated in twice and kickstarted my career and love for radio astronomy. I also want to thank Wael Farah, who I spent many days in lab 339 in Campbell, discussing the world of fast transients and motorcycles. Of course a huge thanks to Dave McMahon and Matt Lebofsky for always helping with backends, password resets, life in the Bay and recommending the local `institutions'. 

When I first got interested in trying my hand in astrophysics as a teenager, I always had wanted to work on a NASA mission. In the Summer of 2024 I was fortunate enough to spend a couple of months developing software for the \textit{Curiosity} rover, I want to thank the SETI institute and the Frontiers Development Lab for making that dream come through. More specifically to the team I worked with Arushi, Kanak and Pugal. To my friends in lab 699 Eric and Victoria, thanks for the hospitality on my visits to the DC area, the tours of Goddard and letting me work on such a cool project. To Robert Alvarez, thanks for being a great mentor and friend, but most importantly showing my what real rum tastes like. 

Thanks to Duncan Lorimer, Maura McLaughlin and all the folks at WVU for hosting me during my visit to GBO and Morgantown. 

Thanks to A-a-ron, who always lent me a desk when I needed it and not abandoning me in the Mongolian steppe when our truck broke down in a snow storm when he really could have gone buckwild on the whole operation. 

Thanks to Hugh for being my gym buddy over the years and to Diarmuid for all his trading and health advice.   

Tiernan, Megan and Ellen thank you guys for always lending an ear when I needed it 

Sam, Ethan and Jonas, thank you guys for all the travels, conversations and enjoyment over the years. 

Oskaras, Sammy and Craig thanks for letting me crash on their couch every time I came to Galway.

I want to thank my Uncle Lawrence, for introducing me to computers and video games as a kid kickstarting my interest in software and hardware alike, along with his endless support and enthusiasm for my efforts in education. 

Thanks to my sister Sinéad, godparents Agnes and Henry for all their support and faith throughout the years. Of course my parents Phyllis and Austin, to which this thesis is dedicated. When I think of all the things I've been through I would have never made it without you guys. 

To Eve, Noel and Bart, I've never forgotten about you. I've tried to carry at least a fraction of your love for the sciences with me throughout this endeavour. Ar dheis Dé go raibh a n-anamacha.

And of course, my partner in crime Laura. Thank you for your unrelenting belief in me over these last four years and always showing up for me on my worst days. This PhD would have been a lot more difficult if it wasn't for you. I am sorry I spent so much time in Birr or at my computer when I really should have been strolling along the canal with you.

\newpage 
\thispagestyle{empty}
\section*{Funding and Resource Acknowledgements}

This PhD was jointly supported by Breakthrough Listen managed by the Breakthrough Prize Foundation and Trinity College Dublin, School of Physics. This work extensively used the Irish LOFAR station, funded by the Irish LOFAR  Consortium. The consortium consists of Trinity College Dublin, Armagh Observatory and  Planetarium, University College Dublin, Dublin City University, University College Cork,  University of Galway, Dublin Institute for Advance Studies and Technological University of  the Shannon Athlone. It receives generous funding from Science Foundation Ireland and the  Department of Further and Higher Education, Research, Innovation and Science.  The REALTA compute cluster was funded by Science Foundation Ireland. This work also used LOFAR-SE 

The work in \cref{sec:Mars} has been enabled by the Frontier Development Lab (FDL4AI.com). FDL is a
partnership between the National Aeronautics and Space Administration (NASA), and the SETI Institute, in collaboration with private industry and academia. This public/private partnership ensures that the latest tools and techniques in Artificial Intelligence (AI) and Machine Learning (ML) are applied to basic science research priorities in support of space science, earth science and energy exploration beneficial to the world. This material is based upon work supported by NASA under Cooperative Agreement No.
NNX14AT27A. Any opinions, findings, and conclusions or recommendations expressed in this material are those of the author(s) and do not necessarily reflect the views of the National Aeronautics and Space Administration.